\documentclass[11pt,a4paper]{article}

\usepackage[margin=1in]{geometry}
\usepackage{parskip}
\usepackage{enumitem}
\usepackage{titlesec}
\usepackage{hyperref}

\titleformat{\section}{\large\bfseries}{}{0em}{}
\titleformat{\subsection}{\normalsize\bfseries}{}{0em}{}

\title{\textbf{Field Visit Report -- Navagamuhena Tea Factory}}
\author{Hasindu Rodrigo, Radeesh Migara, Kalana Ekanayaka}
\date{Visit Date: 2026/06/22}

\begin{document}

\maketitle

\section*{Report Summary}

\textbf{Location:} Navagamuhena Tea Factory

\textbf{Participants:} Sandaru Ruwaneka, Hasindu Nimantha, Radeesh Migara, Kalana Ekanayaka, Santhosh Narayan

\textbf{Prepared By:} Hasindu Rodrigo, Radeesh Migara, Kalana Ekanayaka

A field visit was conducted to Navagamuhena Tea Factory to gain an understanding of the factory's operational processes, existing software solutions, and potential software requirements. The visit provided valuable insights into how technology is currently being utilized within the tea manufacturing industry and helped identify opportunities for the proposed Tea Collection Management System.

\section*{Purpose and Scope}

\subsection*{Purpose}
\begin{itemize}[nosep]
    \item Understand the workflow and operations of a tea factory.
    \item Identify software-related challenges and requirements from the factory's perspective.
    \item Evaluate existing technological solutions used within the industry.
    \item Validate the market demand for a tea collection management solution.
\end{itemize}

\subsection*{Scope}
\begin{itemize}[nosep]
    \item Current software systems used by the factory.
    \item Stakeholder interactions within the tea collection process.
    \item Technological adoption within factory operations, including payment and salary handling.
    \item Potential areas where the proposed system could provide additional value.
\end{itemize}

\section*{Observations}
\begin{itemize}[nosep]
    \item The factory uses a customized software solution to manage several activities, and access is controlled through role-based permissions.
    \item Details of tea collection agents and estate owners are stored in the system, but they do not have direct access to the software.
    \item The factory makes monthly payments to estate owners for supplied tea, after deducting advance amounts and fertilizer costs.
    \item Salary payment details are prepared by the factory accountant and sent to the bank for processing.
    \item The factory keeps a manual error book to record system errors and also stores screenshots or pictures of those errors.
    \item Some data is still collected manually first, such as employee attendance and tea leaf weights, before entering it into the system.
\end{itemize}

\section*{Data Collected}
\begin{itemize}[nosep]
    \item Information about the factory's existing software system and how it is used.
    \item Details of user roles and access control within the system.
    \item Records related to tea collection agents, estate owners, and factory administration.
    \item Information about monthly payments, deductions, and financial handling.
    \item Salary payment process and bank-related documentation.
    \item Error recording method, including manual error book entries and screenshots of issues.
    \item Manual data collection methods for attendance and tea leaf weight records.
\end{itemize}

\section*{Key Findings}
\begin{itemize}[nosep]
    \item The factory already has a digital system, but it is mainly used for internal management and selected stakeholders.
    \item Tea estate owners and collecting agents are included in the records, but they are not fully connected to the system.
    \item Financial activities and salary processing are handled separately from the main operational system.
    \item The factory still depends on manual steps for some activities, showing that the process is not fully automated.
    \item There is a clear need for a more comprehensive system that connects more stakeholders and reduces manual work.
    \item Error tracking is currently manual, which suggests the existing system still needs improvement and better support.
\end{itemize}

\section*{Recommendations and Action Items}
\begin{itemize}[nosep]
    \item Conduct field visits to additional tea factories to broaden the requirements baseline.
    \item Investigate options for direct system access or a simplified interface for estate owners and collecting agents.
    \item Explore digitizing the manual error book into a structured error-logging feature within the proposed system.
    \item Assess opportunities to reduce manual data entry for attendance and tea leaf weight recording, such as through digital weighing or mobile data capture.
    \item Evaluate the feasibility of integrating payment, deduction, and salary processing into a unified platform.
\end{itemize}

\section*{Follow-Ups}
\begin{itemize}[nosep]
    \item Schedule visits to additional tea factories.
    \item Conduct interviews with factory owners and administrators.
    \item Gather detailed requirements from estate owners, collecting agents, and receiving agents.
    \item Analyze existing software solutions used within the industry, including role-based access models.
    \item Refine system requirements based on findings from future visits.
    \item Prepare a comparative analysis of requirements across multiple factories.
\end{itemize}

\section*{Conclusion}

The visit to Navagamuhena Tea Factory confirmed that the factory already relies on a customized, role-based software system for several of its operations, but that estate owners and collecting agents remain outside direct access to it, and key activities such as attendance tracking, tea leaf weighing, and error logging are still handled manually. The proposed Tea Collection Management System offers an opportunity to connect these stakeholders more directly, reduce manual processes, and provide a more structured approach to error tracking and financial handling. Further field visits are necessary to understand varying factory requirements and ensure that the proposed system addresses industry-wide needs effectively.

\end{document}
